% presentation/slides/03_modele.tex
\begin{frame}{Le Modèle Mathématique Général}
    Pour chaque observation $i \in \{1, \dots, n\}$, on formule l'équation :
    
    \begin{center}
        \Large \alert{$Y_i = A_{Z_i}\theta + \sigma E_i$}
    \end{center}
    
    \vspace{0.4cm}
    \begin{itemize}
        \item $\mathbf{\theta \in E}$ : Le paramètre ou signal d'intérêt \textbf{inconnu} à estimer ($\dim E = d$).
        \item $\mathbf{Z_i \in \mathcal{Z}}$ : La \textbf{variable cachée} (la transformation géométrique subie par l'objet), supposée i.i.d. selon une loi connue $f_Z$.
        \item $\mathbf{A_{Z_i}}$ : L'opérateur linéaire modélisant la transformation (translation, rotation, projection).
        \item $\mathbf{E_i \sim \mathcal{N}(0, \text{Id}_p)}$ : Un bruit blanc gaussien i.i.d., indépendant de $Z_i$.
        \item $\mathbf{\sigma > 0}$ : Le niveau de bruit de l'appareil de mesure.
    \end{itemize}
\end{frame}